\documentclass[11pt]{article}
\usepackage[margin=1in]{geometry}
\usepackage{amsmath,amssymb,mathtools,hyperref}
\DeclareMathOperator{\diag}{diag}
\DeclareMathOperator{\rank}{rank}
\title{Native truncation bounds for the original mixed-action singular values}
\author{}
\date{20 September 2026}
\begin{document}\maketitle
This receiving calculation combines the proved native truncation
TR1--10 with DQ1--12 and PM1--20. It keeps the original polynomial
source, its quotient by the full root polynomial, the original
observation and the inherited period Gram. The finite native
integrals in TR are retained as their actual integrals; no moment
evaluation or numerical certificate is asserted here.

\section{The same source bracket through both original minima}
Retain TR's original cutoff $N\ge q-1$, parity-coordinate Gram $H$,
and finite native-resolvent truncation $H^{[J]}$. With the exact
original constant $L_N^\sharp$ from TR3, take the original integer
$J>L_N^\sharp/2$ and put
\[
 \epsilon=\frac{L_N^\sharp}{2J}\in(0,1),\qquad
 \kappa=(1-\epsilon)^{-1}.
 \tag{MS1}
\]
TR1--4 proves $H^{[J]}\preceq H\preceq\kappa H^{[J]}$ from
the complete original measure and multiplication bound. This is an
available result at the displayed integers, not a new assumption.
Let $A_0$ be the original map from parity-coordinate polynomials to
$E_k=\mathbb C[S]/\chi_k$, including the full physical-coordinate
matrix. Its type and kernel are unchanged. Define
\[
 \underline G=(A_0(H^{[J]})^{-1}A_0^*)^{-1},\quad
 G_N=(A_0H^{-1}A_0^*)^{-1},\quad
 \underline Q=(\Lambda\underline G^{-1}\Lambda^*)^{-1}.
 \tag{MS2}
\]
The first onto map is onto because $N\ge q-1$; $\Lambda$ is onto
its actual image. The fixed-fibre minimum formula, proved by
completing the square in TR5 and PM3, gives
\[
 \underline G\preceq G_N\preceq\kappa\underline G,\qquad
 \underline Q\preceq Q_N\preceq\kappa\underline Q,\qquad
 \underline H_K:=I^*\underline GI\preceq H_K
                           \preceq\kappa\underline H_K.
 \tag{MS3}
\]
Here $I$ is the unchanged original frame of $\ker\Lambda$.
Minimizing over the same affine fibre preserves both inequalities:
for every admissible representative the source inequality holds,
and taking infima on that same fibre proves it. This argument
keeps both complete minima in MS2.

\section{Enclosures for every mixed-action singular value}
The original forward mixed-action map is the fixed matrix
$B=\Lambda M_S I$. Its coordinates do not depend on the source
metric. The source and target metric presentations are
\[
 B_N=Q_N^{1/2}BH_K^{-1/2},\qquad
 \underline B=\underline Q^{1/2}B\underline H_K^{-1/2}.
 \tag{MS4}
\]
For every singular-value index of this rectangular map, including
its zero values, the exact native enclosure is
\[
 \boxed{\sqrt{1-\epsilon}\,\sigma_j(\underline B)
       \le\sigma_j(B_N)
       \le(1-\epsilon)^{-1/2}\sigma_j(\underline B).}
 \tag{MS5}
\]
For proof, each nonzero coordinate vector $x$ satisfies
\[
 \kappa^{-1}\frac{x^*B^*\underline QBx}{x^*\underline H_Kx}
 \le\frac{x^*B^*Q_NBx}{x^*H_Kx}
 \le\kappa\frac{x^*B^*\underline QBx}{x^*\underline H_Kx}.
\]
These quotients are the Rayleigh quotients of the two squared
singular-value problems in the same coordinate space. Their
min--max formula follows from an orthonormal eigenbasis after
the indicated metric coordinate isometry: the span of the first
$j$ eigenvectors supplies the lower bound, and intersection with
the orthogonal complement of the first $j-1$ supplies the reverse
bound. Taking the identical extrema over coordinate subspaces
in the two displayed inequalities and then positive square roots
proves MS5. Kernel values remain exactly zero because both metrics
are invertible and the matrix $B$ is unchanged.

The original reverse direction $C=J_KM_SL$ depends on the actual
metric, so it is not replaced by the reverse block formed from
$\underline G$. Instead use the exact proved identity
$C_N+B_N^*=R_N$, $\rank R_N\le2$, where
$C_N=H_K^{1/2}CQ_N^{-1/2}$. Intersecting a $j$-dimensional
singular space of $B_N^*$ with $\ker R_N$ loses at most two
dimensions and on that intersection $C_N=-B_N^*$ exactly.
For every $3\le j\le\min(\dim\ker\Lambda,\dim B_k)$ this proves
\[
 \boxed{\sigma_{j-2}(C_N)
          \ge\sqrt{1-\epsilon}\,\sigma_j(\underline B).}
 \tag{MS6}
\]
The actual $G_N$-metric commutator has the two isometric blocks
$\left(\begin{smallmatrix}0&-C_N\\B_N&0\end{smallmatrix}\right)$.
Its singular values are the union of those of $C_N$ and $B_N$,
since its squared singular matrix is block diagonal. Thus at
least $2j-2$ of its singular values are at least the right side
of MS6 whenever $\sigma_j(\underline B)>0$. This is the exact
quantitative receiver for the original mixed return, without
constructing a false metric-independent reverse block.

\section{Period-norm conversion and a fixed recovery map}
Retain the complete inherited period Gram $P=\jmath^*\jmath$ of
PM4, and put
\[
 \underline\rho=\sqrt{\lambda_{\max}
          (\underline Q^{-1/2}P\underline Q^{-1/2})}.
\]
The sharp original conversion constant in PM5--7 satisfies
\[
 \boxed{\sqrt{1-\epsilon}\,\underline\rho
      \le\rho_N\le\underline\rho.}\tag{MS7}
\]
Indeed $\rho_N^2=\sup_{b\ne0}b^*Pb/(b^*Q_Nb)$; the numerator
is unchanged and MS3 brackets its denominator. Taking the same
supremum gives the result. Hence the actual extra factor required
by PM19 can be enclosed from the same finite native truncation.

For a fixed actual recovery matrix $F:B_k^{\oplus(d+1)}\to E_k$
with $F(\Lambda,\ldots,\Lambda M_S^d)^T=I$, set
$C_N(F)=\|F\|_{Q_N^{\oplus(d+1)}\to G_N}$ and define
$\underline C(F)$ with the two underlined metrics. Applying MS3
to the numerator and denominator of the same operator Rayleigh
quotient gives
\[
 \sqrt{1-\epsilon}\,\underline C(F)
 \le C_N(F)\le(1-\epsilon)^{-1/2}\underline C(F).
 \tag{MS8}
\]
For example the inherited-period recovery of PM8 is such a fixed
map. One may also construct the DO corner left inverses using
$\underline Q$: their exact identity $J_iB_i=I$ holds for every
positive definite output metric, so the resulting reconstruction
remains an exact inverse of the original observation stack.
Once chosen it is held fixed in MS8. This does not equate it
with the different left inverse selected using $Q_N$.

All matrices underlined here still contain the actual native
integrals retained by TR1. Equations MS5--8 propagate certified
values or enclosures of those matrices to the original mixed
directions and the period conversion without removing any kernel
or taking an unproved arithmetic phase limit.

\begin{thebibliography}{9}
\bibitem{TR} Split-Zero programme, \emph{Native resolvent truncation
at growing degree and the original conductor minimum}, TR1--10,
complete proof \texttt{NATIVE\_TRUNCATION\_CONDUCTOR\_RETURN.tex},
published in the edition
\url{https://github.com/KokunoYumeto/zeta-function-research-reader/tree/ab45e221579a0d48ce2885b4ecdf1a6aefc24a20/workbenches/splitzero-tandem/continuations/20260920-native-parity-conductor-return}.
The full proof was read for this receiver. Its human-literature
input is credited to A.~I.~Aptekarev, G.~L\'opez Lagomasino and
A.~Mart\'{\i}nez--Finkelshtein, \emph{Strong asymptotics for the
Pollaczek multiple orthogonal polynomials ensembles},
arXiv:1410.1261v1; the complete companion proves the actual native
transport and reciprocal series used in TR. This note uses those
proved programme maps and does not import the author's ensemble
asymptotics or claim a new reading of the whole author paper.
\end{thebibliography}
\end{document}
